% Chapter 2 — Section: Evolution of Next Best Action Systems and the Research Gap
% Generated from: literature_step1.md + verified_papers.md
% All citations verified in ≥2 databases. Format: \cite{} only.
% Do NOT insert \begin{document} or markdown — paste directly into the chapter.

\subsection{Evolution of Next Best Action Systems and the Research Gap}
\label{subsec:nba-evolution}

The problem of predicting a user's next preferred item — variously framed as session-based
recommendation, sequential prediction, or next best action (NBA) — has evolved through three
distinct algorithmic generations over the past fifteen years.
The earliest formalization by \cite{rendle2010fpmc} combined matrix factorization with
first-order Markov chain transitions, enabling tractable inference over personalized item
sequences even under sparse interaction data.
This Markov foundation established the key inductive bias of the field: that the immediate
prior interaction is the strongest predictor of the next, a principle all subsequent methods
either extend or challenge.
The deep learning transition arrived with \cite{hidasi2016gru4rec}, who demonstrated that
Gated Recurrent Units (GRU) capture session-level sequential dependencies substantially more
effectively than Markov models by modeling unbounded interaction histories through learned
hidden states.
\cite{kang2018self} subsequently replaced recurrence with self-attention, introducing
\emph{SASRec}, which attends selectively to decision-relevant past items via a causally
masked Transformer and achieves superior accuracy on both sparse and dense interaction
datasets.

Despite this rapid architectural progress, traditional sequential recommendation methods share
a fundamental limitation that none of the above systems address: item representations are
learned ID embeddings, carrying no visual attributes, textual descriptions, or multimodal
content signals.
Under this paradigm, two visually identical products with distinct catalog identifiers are
treated as semantically unrelated; models trained on fixed ID vocabularies fail on new or
infrequently purchased items where embeddings remain undertrained; and inference weights are
frozen at serving time, leaving models unable to adapt to within-session preference shifts as
a user's intent evolves.
This shared failure mode — the inability to incorporate item visual attributes, the fragility
of ID-based representations under cold-start conditions, and the rigidity of static weights
under session-level context shifts — constitutes the central research gap that our dual-pipeline
NBA system is designed to address.

The multimodal wave of 2022--2024 began closing this gap by aligning visual and textual item
signals with collaborative filtering.
\cite{wei2023mmssl} demonstrated via inter-modality contrastive learning that visual-textual
alignment provides recommendation signal orthogonal to co-occurrence statistics, substantially
improving cold-start performance on e-commerce benchmarks.
\cite{zhou2023bm3} showed that simple dropout-based contrastive augmentation across modalities
achieves competitive accuracy with 2--9$\times$ lower training cost than graph-based multimodal
methods, establishing visual-textual bootstrapping as a practical industrial approach.
Most recently, \cite{guo2024lgmrec} demonstrated that combining local item-level and global
user-level graph structures with multimodal features consistently outperforms single-modality
ablations, confirming that visual content and interaction structure are complementary signals
rather than substitutes.
Our system operationalizes this multimodal insight through two independent pipelines: Pipeline~A
employs Cleora behavioral co-purchase graph embeddings re-ranked by BGE-M3 content similarity,
targeting the structural signal \cite{rendle2010fpmc} identified; Pipeline~B extends the
\cite{xie2022difsasrec} decoupled side-information architecture with per-user online
fine-tuning, enabling visual and categorical features to modulate sequential intent while
the \texttt{AgentPool} mechanism allows model weights to adapt within a live session.
